% !PRINTED 35
\weberpage{35}
%<<CONT

\begin{equation}
\begin{aligned}
A&=Q'A'+A'',\\
A'&=Q''A''+A''',\\
&\cdots\\
A^{(\nu-3)}&=Q^{(\nu-2)}A^{(\nu-2)}+A^{(\nu-1)},\\
A^{(\nu-2)}&=Q^{(\nu-1)}A^{(\nu-1)}+A^{(\nu)},
\end{aligned}
\tag{3}
\end{equation}
worin $A^{(\nu)}$ eine Constante ist.

Wenn nun $A,A'$ irgend einen gemeinsamen Theiler haben, so ist dieser nach den Sätzen des vorigen Paragraphen, wie der Anblick der Gleichungen (3) lehrt, auch Theiler von $A'',A''',A''''$ u. s. f. bis $A^{(\nu)}$. Ist $A^{(\nu)}$ eine von Null verschiedene Constante, so kann also kein von $x$ abhängiger Theiler von $f(x)$ und $\varphi(x)$ existiren. Solche Functionen heissen \gesperrt{theilerfremd} oder \gesperrt{relativ prim}. Man sagt auch, indem man nur die von $x$ abhängigen Theiler berücksichtigt, die Functionen haben keinen gemeinsamen Theiler.

Die Bedingung, dass die beiden Functionen einen gemeinsamen Theiler haben, ist also:
\begin{equation}
 A^{(\nu)}=0.
\tag{4}
\end{equation}
In diesem Falle ist $A^{(\nu-1)}$ Theiler von $A^{(\nu-2)}$, wie die letzte Gleichung (3) zeigt, und nach der vorletzten dieser Gleichungen Theiler von $A^{(\nu-3)}$ u. s. f., also auch Theiler von $A$ und $A'$, d. h. von $f$ und $\varphi$. Und da umgekehrt jeder gemeinsame Theiler von $A,A'$ auch Theiler von $A^{(\nu-1)}$ ist, so heisst $A^{(\nu-1)}$ der \gesperrt{grösste gemeinsame Theiler} von $f$ und $\varphi$. Der Algorithmus (3) zeigt, dass man $A^{(\nu)}$ oder $A^{(\nu-1)}$ aus den Coëfficienten von $f$ und $\varphi$ durch die \gesperrt{rationalen Rechenoperationen} ableiten kann, und zwar so, dass immer nur durch die Coëfficienten der höchsten Potenzen von $x$ in den Functionen $A,A',A'',\ldots$, die von Null verschieden sind, dividirt wird.

Wir wollen diese Betrachtungen auf zwei Beispiele anwenden.

Es seien zunächst:
\begin{equation}
\begin{aligned}
 A&=a_0x^2+a_1x+a_2,\\
 B&=b_0x^2+b_1x+b_2
\end{aligned}
\tag{5}
\end{equation}
zwei Functionen zweiten Grades, also $a_0,b_0$ von Null verschieden.
